\section{Discussion}
\label{sec:discussion}

\paragraph{The improvement result is narrower than ``PPL is useless.''}
Perplexity remains the direct measure of the continuation objective and is
essential for monitoring CPT. The result is instead about an implication:
after this structural intervention, a lower in-domain PPL does not by itself
imply recovery on MMLU or the tested closed-book tasks. A compressed model may
therefore look healthier under the training-domain likelihood while still
remaining unsuitable for a target use. We did not test a prospective absolute,
relative-to-base, plateau, or calibrated PPL rule.

\paragraph{Null baselines must travel with interface claims.}
The random-init content score changes how the content interface should be read.
Without that null point, keep14's larger content than letter score could be
mistaken for target recovery. With it, the safer conclusion is that content
scoring has a high baseline under this recipe and that interface choice
materially changes the observed score. Neither interface is designated ground
truth; agreement across independently motivated evaluations is more
informative.

\paragraph{Construction labels are part of the result.}
Calling both keep14 and ShortGPT ``16-layer models'' hides inherited count,
selected indices, final-block retention, and fresh-tail use. Their different
endpoints show why exact construction must accompany quality claims. The
present comparison cannot say which difference matters, and random/frozen
likewise cannot isolate initialization or adaptation. Reporting these rows as
operating points preserves their falsification value without converting them
into causal ablations.

\paragraph{Recommended reporting axes.}
For post-intervention recovery, we recommend reporting:
(1) in-domain and, when available, out-of-domain likelihood;
(2) the target evaluation and its sample count;
(3) prompt/scoring/generation interface and a relevant null baseline;
(4) exact inherited/fresh structure and trainable modules;
(5) steps, token presentations, realized compute, and stopping rule; and
(6) independent training runs separately from item-level uncertainty.
This is a case-study recommendation, not an empirically validated universal
standard and not a substitute for latency, memory, or endpoint-quality
frontiers.
